\section{Threat-model boundary: what consolidation actually removed}\label{sec:boundary}

The controlled experiments above install the corruption by hand, deleting a
negative caveat entirely. Whether that omission arises on its own is a separate
question, and the answer bounds everything else in this paper.

\paragraph{Four things a memory can lose.} These come apart in what follows and
are not interchangeable. A \textbf{negative outcome} is an adverse empirical
result, here retention $-12$pp. A \textbf{scope constraint} restricts the
population, period or context a result holds for. A \textbf{prohibition} is a
normative instruction, ``do not reuse this''. \textbf{Complete negative
deletion} is the corruption our controlled sections install: no adverse outcome
and no warning of any kind remains.

Figure~\ref{fig:erosion} summarises what happened. We ran a benign consolidation pipeline --- \FiveChains{} chains,
\FiveGenerations{} generations, each generation asked only to compress and carry
forward, with no instruction about caveats or qualifiers --- and scored every
generation of every chain with a frozen model scorer. On the two records that
carry no qualifier in their source it returns \FiveFalsePosAll{} false
positives, and on the four that do it recovers the qualifier in
\FiveScorerSensK/\FiveScorerSensN{} first-generation bodies. Two limitations
follow it into every number below: the scorer is \texttt{\FiveScorerModel},
which is \textbf{also one of the six consolidator models}, so it is not
independent of the systems it judges; and no human validation was performed.
Its sensitivity is validated only on first-generation bodies, which are
near-verbatim; nothing validates it on the compressed final-generation bodies
where the erosion is claimed, so part of the decline could be scorer failure on
telegraphic text rather than loss of the qualifier. The negative-outcome
counts, by contrast, are deterministic keyword detection and do not depend on
the scorer at all --- and they are immune to this confound.

\paragraph{What survived.} \textbf{Complete negative deletion} appears in
\FivePurelyPositive{} chain-generations. This is the only claim the
\FivePurelyPositive{} figure supports; it is not a claim that nothing eroded. The quantified negative is durable: it is present in
\FiveNumAllGens{} chains at \emph{every} one of the six generations
(\FiveNumSeries\%), and of the \NineBodies{} distinct final-generation bodies,
\NineKeepNeg{} carry a number attached to a negative outcome or a churn or
retention term. Appendix~\ref{app:consolidation} prints final-generation bodies
verbatim, selected by a stated deterministic rule.

\paragraph{What did not.} The same scorer shows the qualifier eroding
\emph{monotonically inside the horizon we tested}. For the target memory the
share of chains still stating its qualifier runs
\FiveTgtSeries\% across the six generations, ending at
\FiveTgtGenSixK/\FiveTgtGenSixN. For the three memories whose qualifier is a scope
restriction rather than a number, it runs \FiveScopeSeries\%, ending at
\FiveScopeGenSixK/\FiveScopeGenSixN. The explicit prohibition falls too, from
\FiveProGenOneK/\FiveProGenN{} chains to \FiveProGenSixK/\FiveProGenN{}
(\FiveProSeries\%), though not monotonically, so we report the level rather
than a trend. Consolidation did not remove the number; the scope
declined monotonically and roughly half the prohibitions were gone by the sixth
generation.

\paragraph{What that does to the risky action.} Running the allocation protocol
on the final-generation bodies, a consolidated body suppresses the risky action
as completely as the intact original: it is planned in
\NineNatIntentK/\NineNatIntentN{} episodes against
\NineIntactIntentK/\NineIntactIntentN{} for the intact original and
\NineHandIntentK/\NineHandIntentN{} for the hand-stripped body. It is verified far less often than our hand-built corruption. The
prospectively specified primary contrast --- the consolidated body padded to the
hand-built body's length, because \S\ref{sec:content} shows a
\EightLengthEffect-point silent-deletion versus neutral-clause gap --- is
\NinePadVsHand{} points; unpadded, the raw gap is \NineNatVsHand{} points
(\NineNatPct\% against \NineHandPct\%). Against the intact original it is
\NineNatVsIntact{} points, a difference we do not interpret: it is smaller than
the run-to-run variability we measure in \S\ref{sec:limits} and has not been
replicated (full arm table in Appendix~\ref{app:models}). \textbf{We cannot say whether a consolidated body is harder or
easier to catch than the record it came from}; what we can say is that it is
not the corruption our controlled sections install.

\begin{center}\itshape
The allocation dependence is real, but whether it becomes harmful depends on
what memory degradation actually removes. In the benign consolidation process we
tested, the quantified negative survived all six generations and complete deletion
did not occur --- while scope fell to \FiveScopeGenSixPct\% and roughly half
the explicit prohibitions were gone.
\end{center}

\noindent We do not generalise this beyond the tested scenario, models,
consolidation procedure and \FiveGenerations{}-generation horizon, and the
erosion trend is a reason not to: it is a decline observed \emph{within} the
horizon, not a plateau. A longer horizon, a lossier or adversarial pipeline, or
a constraint expressed only as scope may well produce the omission this pipeline
did not.
